A (simple) graph $G$ has chromatic number $\chi(G)$ equal to the least cardinal $\kappa$ such that $V(G)$ can be 
properly colored with $\kappa$ colors.

Given graphs $G_1,G_2$ with $\chi(G_1)=\chi(G_2)=\aleph_1$,
question: must there exist a graph $H$ with $\chi(H)=4$ (or alternatively with $\chi(H)=\aleph_0$) such that $H$ is a 
subgraph of both $G_1$ and $G_2$? Here ``$H$ is a subgraph of $G_i$'' means there exists an injective map 
$f_i:V(H)\to V(G_i)$ such that if $xy\in E(H)$ then $f_i(x)f_i(y)\in E(G_i)$.
